\documentclass[11pt]{letter}
\usepackage[margin=1in]{geometry}
\usepackage[utf8]{inputenc}
\usepackage[T1]{fontenc}
\usepackage{microtype}
\usepackage{hyperref}

\hypersetup{colorlinks=true, urlcolor=blue!70!black}

\signature{Drake Caraker\\Oura Health, San Francisco, CA\\
\texttt{drakecaraker@gmail.com}}

\address{Drake Caraker\\Oura Health\\San Francisco, CA}

\begin{document}

\begin{letter}{The Editors\\Nature}

\opening{Dear Editors,}

We submit for your consideration an Article entitled ``The Limits of
Explanation,'' which proves a universal impossibility theorem for
explanation of underspecified systems, unifying known results across
eight scientific fields under a single structural cause.

\textbf{The result.}  We prove that no explanation of an
underspecified system can simultaneously be faithful (consistent with
the system's structure), stable (invariant across observationally
equivalent configurations), and decisive (committing to every
structural distinction).  This impossibility is not a conjecture or an
empirical observation: it is a theorem, derived from first principles
in each of eight domains---physics (gauge freedom), biology (codon
degeneracy), statistics (Markov equivalence), crystallography (the
phase problem), linguistics (syntactic ambiguity), mathematics
(underdetermined linear systems), computer science (the view-update
problem), and statistical mechanics (microstate multiplicity).  These
phenomena were previously studied in complete isolation.  We show they
are instances of a single impossibility, and that the resolution each
field independently adopted---averaging over equivalent
configurations---is Pareto-optimal.

\textbf{Empirical validation.}  Three quantitative predictions of the
framework have been confirmed: invariance counting in codon degeneracy,
the Gaussian flip rate for feature attributions, and the universal
$\eta$~law governing explanation instability.  Two pre-registered
extensions were falsified (a predicted phase transition in Rashomon set
volume and a conjectured uncertainty bound), demonstrating the
framework's testable boundaries.  The constructive resolution
($G$-invariant projection) is proved Pareto-optimal.  The entire
formal framework is mechanically verified in the Lean~4 proof
assistant (491~theorems, 0~unproved goals).

\textbf{Why Nature.}  The result is inherently cross-disciplinary: a
biologist, a physicist, and a computer scientist can each follow the
paper because each sees their own domain's impossibility derived from
shared structure, not borrowed formalism.  This cross-disciplinary
unification is Nature's distinctive strength and would not be served by
a single-domain venue.

\textbf{Practical significance.}  The EU AI Act (Article~86) requires
transparent explanations of automated decisions.  This paper proves
that the demanded trifecta---faithful, stable, and decisive
explanations---is mathematically impossible under underspecification,
and provides the optimal achievable alternative.

\textbf{Companion papers.}  A companion paper on the physics
application has been submitted to \textit{Foundations of Physics}, and
an ML-specific paper providing domain depth for feature attributions is
in preparation.  This submission focuses on the unification.

\textbf{Suggested reviewers.}  The paper touches ML interpretability,
impossibility theorems, and formal verification.  We suggest
interdisciplinary reviewers spanning (i)~formal methods or proof
assistants, (ii)~causal inference or statistical methodology, and
(iii)~physics or biology, to evaluate the cross-domain derivations.

We confirm that this work is original, has not been published
elsewhere, and is not under consideration at any other journal.  All
authors have approved the manuscript.  We declare no competing
interests.

\closing{Sincerely,\\[0.5em]
Drake Caraker, Bryan Arnold, David Rhoads}

\end{letter}
\end{document}
